\section{System and Threat Model}
\label{sec:threat-model}

\subsection{System model}
\label{sec:system-model}

We consider a deployment of $N=24$ nodes: eight sensors, eight actuators and
eight mixed nodes that both sense and actuate. Nodes communicate through an edge
gateway that hosts the policy decision point (PDP); each node hosts a policy
enforcement point (PEP) that executes the selected action or denies the request.
The cloud is used for audit and model updates but is not in the control loop.
Three hardware classes are modelled: a Class-1 microcontroller (Cortex-M4-class,
64--256\,KB SRAM, 10--30\,mW), a Class-2 IoT node (ESP32-S3-class, 512\,KB SRAM,
120--350\,mW) and an edge gateway (Raspberry-Pi-class, multi-watt). The tier mix
is 55\,\% Class-1, 30\,\% Class-2 and 15\,\% gateway.

Each epoch $t$ produces a context observation $\bx_t\in\mathbb{R}^9$ with nine
features: battery level, CPU load, temperature, log packet rate, packet loss,
log round-trip time, an anomaly score, failed-authentication count, and
cross-node dispersion (the interquartile range of a feature across neighbouring
reports). The features are grouped into four modalities
$\mathcal{M}=\{\mathrm{dev},\mathrm{net},\mathrm{sec},\mathrm{cross}\}$:
device telemetry (battery, CPU, temperature), network telemetry (rate, loss,
RTT), security telemetry (anomaly, failed authentication) and cross-node
dispersion. Device telemetry is the \emph{attestation anchor}: it is produced by
the node's trusted execution environment and is assumed to be the least
spoofable modality. Cross-node dispersion is always retained at full weight
because it is the only feature that directly exposes a colluding minority.

A request has an operation class
$\op\in\{\mathrm{read},\mathrm{write},\mathrm{actuate},\mathrm{safety}\}$.
The action space is the Cartesian product of five mechanism dimensions:
authentication $\{$\emph{token}, \emph{mutual}, \emph{attestation}$\}$,
cryptography $\{$\emph{none}, \emph{AES-GCM}, \emph{Ascon}$\}$,
inspection $\{$\emph{off}, \emph{header}, \emph{deep}$\}$,
rate limiting $\{$\emph{unlimited}, \emph{50/s}, \emph{10/s}$\}$ and freshness
$\{$\emph{300\,s}, \emph{60\,s}, \emph{10\,s}$\}$. The lattice therefore has
$3^5=243$ actions plus a deny action. Each action has a latency and an energy
cost (analytical model, \cref{sec:implementation}) and a detection probability
against a malicious request that depends on the latent threat state.

\subsection{Threat model}
\label{sec:threat-model-details}

\paragraph{Adversary goals.} The adversary wants to (i) cause a physical action
that the legitimate operator would not authorise, (ii) prevent a legitimate
safety action, or (iii) exhaust the deployment's energy or availability by
forcing unnecessary security work. Goals (i) and (ii) are safety failures; goal
(iii) is a denial-of-sleep / economic failure.

\paragraph{Capabilities.} The network adversary is Dolev--Yao: it can eavesdrop,
inject, replay, reorder and drop messages. It has physical access to a minority
of nodes and compromises $f=10$ of 24 ($f< N/2$), which then report fabricated
context and may collude (Byzantine). It can poison context in both directions:
spoof a benign context to induce a downgrade, or spoof threat indicators to
induce over-reaction. It can pulse at the adaptation period to induce thrashing,
and it can adapt its strategy over time. The adversary is white-box: it knows
the policy, thresholds, invariant set, action lattice and simulator. It does not
know future random draws and is not assumed to solve the defender's POMDP
optimally; the adaptive scenario is a stress test, not an optimality claim.

\paragraph{Trust assumptions.} The gateway and its PDP are trusted and
attested; secure boot protects the PEP; the certificate authority and supply
chain are not compromised; the device attestation anchor is available on each
node. If a node's own trusted execution environment is subverted, I1 still
limits the damage that node can do to \emph{other} physical processes, but it
cannot prevent the node from exercising credentials it legitimately holds. The
\texttt{tee} scenario stresses exactly this out-of-model condition; it is
reported separately and is not used to claim that the architecture detects a
TEE compromise. It also instantiates the fifth latent state, \texttt{compromised},
which is not generated by the other scenarios.

\paragraph{Out of scope.} Physical destruction of the gateway, side-channel key
extraction, compromised certificate authority, supply-chain implants and
wide-area network partition are out of scope. So are human factors and
insider misuse by an operator with valid credentials; the model addresses
context manipulation, not authorised intent.

\subsection{Hard safety invariants}
\label{sec:invariants}

The defender commits to four invariants. They are properties of the enforcement
point, not of the estimator, and they hold for every belief---including a belief
chosen by the adversary.

\begin{description}
  \item[I1 (actuation envelope).] Every \emph{actuate} action that is granted
  uses mutual or attestation authentication, an AEAD construction, freshness of
  at most 60\,s, and at least header-level inspection. No battery, load or
  latency condition may select plaintext, unauthenticated, stale or uninspected
  actuation.
  \item[I2 (safety envelope).] Every \emph{safety}-critical action additionally
  uses deep inspection and mutual authentication. Low power may reduce telemetry
  sampling; it may not reduce inspection depth or session freshness.
  \item[I3 (distrust restricts).] If the belief mass on serious attack states
  exceeds $\tau_{\mathrm{risk}}$, non-safety operations must use at least header
  inspection and rate limiting; above $\tau_{\mathrm{high}}$ they must also use
  mutual authentication and AEAD. I3 narrows the admissible set; it never
  widens it.
  \item[I4 (per-operation failure semantics).] Deny is admissible for read,
  write and actuate, but not for safety-critical operations. On controller
  timeout, invariant-check failure or an empty admissible set, the PEP executes
  the precomputed safe action for that operation class: fail-closed for
  confidentiality and integrity operations, and fail-safe-allow (authenticated
  emergency function only) for safety-critical ones.
\end{description}

I1 and I2 are \emph{unconditional}: they do not depend on $\attack_t$. I3 is
\emph{conditional}: it depends on the belief and can only remove actions. I4
defines what happens when the controller cannot decide. The separation matters:
a spoof can change I3, but it cannot change I1/I2, and it cannot make a safety
operation deniable.

\subsection{Design goals}
\label{sec:goals}

The design is evaluated against five goals, one per research question.

\begin{description}
  \item[G1 (safety).] No granted \emph{actuate} or \emph{safety} request may
  use an action outside the I1/I2 floor, under any of the eight scenarios.
  Read and write have no structural floor by design, so the IVR is vacuous for
  those classes; we report it separately. Measured as invariant violations per
  1{,}000 decisions (IVR).
  \item[G2 (utility).] Adaptation should retain most of the mission utility of a
  static-high posture and beat static-low, threshold and unconstrained
  risk-adaptive baselines. Measured as normalised utility, false-denial rate and
  residual malicious-grant rate.
  \item[G3 (belief quality).] The risk estimate should be calibrated and should
  detect context poisoning when the evidence permits it. Measured as Brier
  score, ECE, poisoning AUROC and TPR at 1\,\% false-positive rate.
  \item[G4 (fail-safe semantics).] Safety-critical requests should almost never
  be denied by policy, and never by the invariant layer. Measured as the
  safety-denial rate and by exhaustive enumeration of the action lattice.
  \item[G5 (stability).] Adaptation should not thrash and should fit the hardware
  envelope. Measured as switches per 1{,}000 decisions,
  p50/p99 latency, energy per decision and footprint.
\end{description}
